% ============================================================================
%  E — file phần: Hành nghề và luyện tập
%  Ngày tạo: 2026-09-06 | Sửa gần nhất: 2026-09-06 | Ôn gần nhất: chưa
% ============================================================================
\documentclass[../algorithm.tex]{subfiles}
\begin{document}
\part{HÀNH NGHỀ VÀ LUYỆN TẬP}
\begin{tomtat}
Phần E là khoảng cách giữa ``tôi biết phải làm gì'' và ``chương trình chạy đúng, đúng hạn''. Chương 25 nói về cài đặt, gỡ lỗi và kiểm thử --- trong đó stress test là kỹ thuật đáng giá nhất mà rất nhiều người học muộn. Chương 26 nói về thuật toán trong hệ thống thật, nơi hằng số, bộ nhớ và tính bảo trì thường thắng độ phức tạp tiệm cận. Chương 27 là chiến thuật phỏng vấn và thi đấu: quản lý thời gian, chọn bài, nói ra suy nghĩ. Chương 28 là lộ trình luyện tập và ngân hàng bài tập lớn để mang theo sau khi gấp sách. Nếu chỉ nhớ một điều: năng lực giải thuật tăng theo số bài bạn \emph{giải xong sau khi đã bí}, không theo số bài bạn đọc lời giải.
\end{tomtat}

Phần này tồn tại vì một quan sát khó chịu: rất nhiều người hiểu thuật toán tốt hơn hẳn kết quả họ đạt được. Nguyên nhân gần như luôn nằm ở bốn chỗ, và bốn chương của phần E ứng đúng với bốn chỗ đó. Thứ nhất, code có bug mà không có quy trình tìm bug --- chương 25. Thứ hai, chọn sai thứ để tối ưu trong bối cảnh thật, ví dụ đổi cấu trúc dữ liệu để giảm độ phức tạp trong khi nút cổ chai nằm ở truy cập bộ nhớ --- chương 26. Thứ ba, chiến thuật kém: dành 50 phút cho bài khó nhất rồi không kịp làm hai bài dễ --- chương 27. Thứ tư, luyện tập không có cấu trúc: giải nhiều mà không tiến, vì luôn giải những bài vừa sức thay vì những bài hơi quá sức --- chương 28.

Ba chương đầu ngắn và thực dụng; chương 28 dài vì nó chứa ngân hàng bài tập được phân loại theo chương, kèm gợi ý và lộ trình theo tuần. Đó là chương bạn sẽ mở lại nhiều lần nhất sau khi đã đọc xong cả quyển.
\subfile{00-map}
\subfile{25-cai-dat-go-loi-kiem-thu/cai-dat-go-loi-kiem-thu}
\subfile{26-thuat-toan-trong-he-thong-that/thuat-toan-trong-he-thong-that}
\subfile{27-phong-van-va-thi-dau/phong-van-va-thi-dau}
\subfile{28-lo-trinh-va-ngan-hang-bai-tap/lo-trinh-va-ngan-hang-bai-tap}
\ifSubfilesClassLoaded{\printbibliography[title={Nguồn}]}{}
\end{document}
